\section{Classification des formes d'intersection}
%========================================================================

\subsection{Formes unimodulaires indéfinies}

On rappelle que la forme d'intersection $Q_X$ d'une $4$-variété compacte orientée sans bord est unimodulaire, c'est-à-dire que sa matrice dans n'importe quelle base de $H_2(X;\mathbb{Z})$ est de déterminant $\pm 1$. La question naturelle est alors~: quelles formes unimodulaires peuvent apparaître, et comment les distinguer ? Dans le cas indéfini, la réponse est donnée par le théorème de Serre, dont l'énoncé est remarquablement simple.

\begin{theorem}[Serre]
  Deux formes bilinéaires symétriques unimodulaires indéfinies sur $\mathbb{Z}$ sont isomorphes si et seulement si elles ont même rang, même signature et même parité.
\end{theorem}

Autrement dit, les trois invariants introduits à la section précédente suffisent à classifier entièrement les formes unimodulaires indéfinies. En particulier, il n'y a qu'un nombre fini de telles formes pour chaque valeur du rang.

On peut rendre cette classification encore plus explicite. Dans le cas impair, les seules formes possibles sont les sommes directes de copies de $\langle 1 \rangle$ et $\langle -1 \rangle$. Dans le cas pair, interviennent deux formes de base~: la forme hyperbolique
\[
  H = \begin{pmatrix} 0 & 1 \\ 1 & 0 \end{pmatrix}
\]
de rang $2$ et de signature nulle, et la forme $E_8$, définie ci-dessous, qui est de rang $8$ et de signature $\pm 8$.
\[
  E_8 =
  \begin{pmatrix}
     2 & -1 &  0 &  0 &  0 &  0 &  0 &  0 \\
    -1 &  2 & -1 &  0 &  0 &  0 &  0 &  0 \\
     0 & -1 &  2 & -1 &  0 &  0 &  0 & -1 \\
     0 &  0 & -1 &  2 & -1 &  0 &  0 &  0 \\
     0 &  0 &  0 & -1 &  2 & -1 &  0 &  0 \\
     0 &  0 &  0 &  0 & -1 &  2 & -1 &  0 \\
     0 &  0 &  0 &  0 &  0 & -1 &  2 &  0 \\
     0 &  0 & -1 &  0 &  0 &  0 &  0 &  2 \\
  \end{pmatrix}.
\]

\begin{corollary}
  Toute forme bilinéaire symétrique unimodulaire indéfinie sur $\mathbb{Z}$ est isomorphe à l'une des formes suivantes~:
  \[
    m\langle 1 \rangle \oplus n\langle -1 \rangle, \qquad m,n \geq 1 \quad \text{(cas impair)}
  \]
  \[
    kE_8 \oplus \ell H, \qquad k \in \mathbb{Z},\, \ell \geq 1 \quad \text{(cas pair)}
  \]
  où dans le second cas, $k > 0$ correspond à des copies de $E_8$ et $k < 0$ à des copies de $-E_8$.
\end{corollary}

En particulier, la forme d'intersection d'une $4$-variété compacte orientée sans bord dont la forme est indéfinie est nécessairement de l'une de ces deux formes.

\subsection{Le cas défini : une classification impossible}

Contrairement au cas indéfini, la classification des formes unimodulaires définies positives est un problème ouvert en général, et ce dès des rangs relativement petits. En rangs $1$ à $7$, la seule forme unimodulaire définie positive est $\langle 1 \rangle^{\oplus r}$, ce qui reste très simple. Mais en rang $8$ apparaît pour la première fois la forme $E_8$, et à partir de là le nombre de classes d'isomorphisme croît de façon vertigineuse : en rang $16$ il y en a déjà $2$, en rang $24$ plus de $10^9$, et au-delà toute classification exhaustive est hors de portée.

Cette situation contraste fortement avec le cas indéfini, et rend la question de savoir quelles formes sont effectivement réalisées comme forme d'intersection d'une $4$-variété d'autant plus délicate. C'est précisément là qu'interviennent les théorèmes de Freedman et de Donaldson, qui apportent des réponses partielles mais profondes.
